\lab{3}{Токенизация}
\task
Разработать правила выделения токенов и реализацию для дальнейшей
индексации. Описать достоинства, недостатки и неудачные случаи. Измерить
число токенов, среднюю длину, время обработки и скорость на килобайт;
обсудить зависимость от размера данных и способы ускорения.

\method
Токенизатор выделяет последовательности букв Unicode, десятичных цифр
и подчёркиваний. Комбинируемые знаки допускаются после начала сегмента.
Последовательность из одних подчёркиваний не выдаётся. Регистр понижается
по таблице Unicode 15.0; полное case folding и NFC-нормализация не
выполняются. Порядок и повторы сохраняются, стоп-слова не удаляются.

Для технического корпуса введены два дополнения. Разделитель
\code{::} сохраняется между корректными сегментами, поэтому
\code{std::vector::iterator} остаётся одним термином. Обозначения
\code{C++}, \code{G++}, \code{Clang++}, в том числе с номером версии,
также сохраняются. Обычное увеличение \code{value++} таким правилом
не защищается. Это компромисс между естественным текстом и исходным кодом,
а не полноценный синтаксический анализ C++.

\begin{table}[H]\centering
\begin{tabularx}{\textwidth}{lX}\toprule
Вход & Результат\\\midrule
\code{Memory allocation} & \code{memory}, \code{allocation}\\
\code{pthread_create} & \code{pthread_create}\\
\code{std::vector<int>} & \code{std::vector}, \code{int}\\
\code{C++20 value++} & \code{c++20}, \code{value}\\
\code{memory-mapped} & \code{memory}, \code{mapped}\\
\code{can't -O2 6.12} & \code{can}, \code{t}, \code{o2}, \code{6}, \code{12}\\
\code{MemoryBarrier} & \code{memorybarrier}\\\bottomrule
\end{tabularx}
\caption{Примеры правил и их ограничений.}
\end{table}

Строгий UTF-8-декодер отвергает обрезанные последовательности,
недопустимые кодовые точки и ошибочные байты. Длина токена считается
в кодовых точках, включая сохранённые символы \code{_}, \code{:},
\code{+}. Это не длина UTF-8 в байтах и не число визуальных графем.
Среднее вычисляется по всем вхождениям:
\[
\overline{\ell}_{token}=\frac{\sum_{i=1}^{T}\ell(t_i)}{T}.
\]
Библиотечная функция вызывает callback для каждого токена. Самостоятельные
CLI и обработчик \code{corpus.bin} используют один и тот же код.

\journal
Базовые правила для текста дополнены сохранением имён API и обозначений
языка. Затем добавлены таблицы Unicode, строгая проверка входа и
статистический режим без печати каждого токена. Измеритель потока отдельно
считает время внутри вызовов токенизатора и полное время чтения корпуса.
При включённых частотах callback обновляет словарь, поэтому такой замер
нельзя напрямую сравнивать с запуском без сбора частот.

\testing
Табличные C++-тесты покрывают пустой ввод, пробелы, пунктуацию, повторы,
стоп-слова, идентификаторы, пространства имён, версии, пути, числа,
кириллицу, диакритику, комбинируемые знаки и символы вне BMP.
Проверяются не только строки, но и длины. Отдельные случаи неверного
UTF-8 должны завершаться исключением. Python-тесты проверяют извлечение,
профили сайтов, экспорт и чтение корпуса. Таким образом, токенизация и
качество предварительного извлечения проверяются раздельно.

\results
\begin{table}[H]\centering
\begin{tabular}{lrr}\toprule
Показатель & Префикс 10\,000 & Весь корпус\\\midrule
Документов & 10\,000 & 1\,679\,971\\
Текст, байт & 76\,863\,027 & 17\,236\,388\,223\\
Токенов & 10\,683\,808 & 2\,483\,697\,508\\
Средняя длина, кодовых точек & 5,5300 & 4,8031\\
Время токенизации, с & 3,073 & 401,114\\
Полное время, с & 3,392 & 434,245\\
Скорость, MiB/с & 23,853 & 40,981\\
\bottomrule\end{tabular}
\caption{Повторные измерения токенизации. В префиксе включён сбор частот; на полном корпусе он отключён.}
\end{table}
Для полного корпуса время на KiB равно 23,830 мкс. Время
токенизации не включает чтение и расчёт отпечатка, полное время включает их.
Поле \code{terms=0} полного запуска означает отключённый сбор словаря,
а не отсутствие уникальных термов.

В основном индексе сохранено 2\,483\,697\,508 позиций токенов:
эта величина относится ко всему корпусу и совпадает у булевой и
координатной индексации. Она больше числа слов экспортёра, поскольку,
например, дефис, точка и некоторые элементы программного текста
создают дополнительные границы.

По длине входа основной проход линейный. Поиск свойств символа использует
бинарный поиск по фиксированным таблицам Unicode; его стоимость ограничена
размером таблиц, а не корпуса. Для проверки масштабирования можно
использовать несколько префиксов одного и того же потока, сохраняя
одинаковые настройки сбора частот. Скорость первого префикса не следует
считать скоростью случайного среднего документа: в начале находятся
преимущественно страницы документации.

\conclusions
Получен общий токенизатор для корпуса и запросов. Сохранение технических
обозначений уменьшает потерю смысла, но \code{std::vector} и
\code{vector} остаются разными словами, а апострофы, пути и ключи
командной строки обрабатываются грубо. Возможны дополнительные формы
идентификатора, разбиение CamelCase и NFC-нормализация; они потребуют
согласованного изменения индекса и запросов. Ускорить обработку можно
быстрым путём для ASCII и параллельной обработкой независимых документов.
